\clearpage
\begin{summarybox}
\textbf{C086\quad 双曲线焦距比与离心率取值范围——快记／快懂／快用}

\textbf{快记（一个存在性判据）：}
\[
\boxed{\exists P\in H,\ PF_1=\lambda PF_2}
\quad\Longleftrightarrow\quad
\boxed{\lambda>0,\ \lambda\ne1,\quad
1<e\le\frac{1+\lambda}{|1-\lambda|}}.
\]
给定题设已知 $\lambda>0$ 时，需额外排除 $\lambda=1$。

\textbf{快懂（最低记忆负担：差、比、和）：}
\[
\underbrace{|PF_1-PF_2|=2a}_{\text{差固定}},\qquad
\underbrace{PF_1:PF_2=\lambda:1}_{\text{比例已知}},\qquad
\underbrace{PF_1+PF_2\ge2c}_{\text{和不小于底边}}.
\]
所以
\[
\frac ca\le\frac{PF_1+PF_2}{|PF_1-PF_2|}
=\frac{\lambda+1}{|\lambda-1|}.
\]
\textbf{等号：}点 $P$ 在\textbf{对应支顶点}。$\lambda\to1$ 不能替代 $\lambda=1$。

\textbf{快用（识别信号与直接套用）：}
\begin{itemize}
\item 关键词：\textbf{存在点＋两焦距比＋离心率范围}，优先 C086；
\item $PF_1=2PF_2$：$1<e\le3$（右支）；
\item $PF_1=\frac12PF_2$：$1<e\le3$（左支）；
\item $e$ 已定：右支 $1<\lambda\le(e+1)/(e-1)$；
左支 $(e-1)/(e+1)\le\lambda<1$；
\item 需要求 $P$ 的坐标时，再用焦半径公式求 $x$、回代双曲线求 $y$。
\end{itemize}

\textbf{证明方法选择：}
\textbf{最快}——定义＋三角形不等式直接作比；
\textbf{最低门槛}——两腰 $\lambda t:t$，差为 $2a$，和至少为 $2c$；
\textbf{常规}——分左右支，用焦半径公式＋$|x|\ge a$；
\textbf{几何升级}——阿波罗尼斯圆：\emph{定比成圆，定差成双曲线}，
通过轨迹相交判断存在性与切点；
\textbf{构造补充}——以两焦点为圆心的\emph{两个辅助圆}，证明充分性并数交点。

\textbf{阿波罗尼斯圆速记（$\lambda\ne1$）：}
\[
\boxed{C_\lambda\left(\frac{c(\lambda^2+1)}{\lambda^2-1},0\right),\qquad
R_\lambda=\frac{2c\lambda}{|\lambda^2-1|}}.
\]
$\lambda>1$ 圆在右侧，$0<\lambda<1$ 圆在左侧；$\lambda=1$ 时平面轨迹为 $x=0$。
\textbf{圆上极值推理：}对 $\lambda>1$，
$PF_2\in[2c/(\lambda+1),\,2c/(\lambda-1)]$，
距离差最小为 $2c(\lambda-1)/(\lambda+1)$，
从而要求 $2a$ 不小于此数。

\textbf{防错三件事：}\textbf{ $\lambda\ne1$；上界等号能取；$e=1$ 不能取。}

\textbf{可迁移思想：}\emph{差定＋比定 $\Rightarrow$ 和定；两段能否组成三角形决定存在性。}
\end{summarybox}
